\documentclass{article}
\usepackage[T2A]{fontenc}
\usepackage[utf8]{inputenc}   
\usepackage[english, russian]{babel}

% Set page size and margins
% Replace `letterpaper' with`a4paper' for UK/EU standard size
\usepackage[a4paper,top=2cm,bottom=2cm,left=2cm,right=2cm,marginparwidth=1.75cm]{geometry}

\usepackage{amsmath}
\usepackage{graphicx}
\usepackage[colorlinks=true, allcolors=blue]{hyperref}
\usepackage{amsfonts}
\usepackage{amssymb}
% \usepackage[left=1cm,right=1cm,top=1cm,bottom=1cm]{geometry}
\usepackage{hyperref}
\usepackage{seqsplit}
\usepackage[dvipsnames]{xcolor}
\usepackage{enumitem}
\usepackage{algorithm}
\usepackage{algpseudocode}
\usepackage{algorithmicx}
\usepackage{mathalfa}
\usepackage{mathrsfs}
\usepackage{dsfont}
\usepackage{caption,subcaption}
\usepackage{wrapfig}
\usepackage[stable]{footmisc}
\usepackage{indentfirst}
\usepackage{rotating}
\usepackage{pdflscape}
\usepackage{listings}

\usepackage{minted}

\usepackage{MnSymbol,wasysym}

\title{Домашнее задание 3}
\author{Лиза Фоменко}
\date{}

\begin{document}

\maketitle



\section*{\textbf{Задание 1}} \textbf{\textit{Злодей Анти-человек придумал последовательность чисел Анти-начи. Она продолжает последовательность чисел Фибоначчи влево. Например, поскольку $F_0 = 0$, $F_1 = 1$ и $F_{k+2} = F_{k+1} + F_k$, выполняется равенство $1 = 0 + F_{-1}$, из чего следует, что $F_{-1} = 1$. Дальнейшие числа Анти-начи определяются аналогично.
Анти-человек умеет быстро возводить матрицы в степень. Подскажите, как ему находить $F_k$ для отрицательных $k$.}}.

Поступим с анти-начи так же, как с числами Фибоначчи. 

Мы значем, что: 
\begin{enumerate}
    \item $F_0 = 0$
    \item $F_{-1} = F_1 - F_0 = 1 - 0 = 1$
    \item $F_{-2} = F_0 - F_{-1} = 0 - 1 = -1$
    \item $F_{k} = F_{k+2} - F_{k+1}$, где где $k \in \mathbb{N}$ (рекуррентная формула для анти-начи)
    \item $F_{-k} = F_{-k+2} - F_{-k+1}$
\end{enumerate}

Давайте теперь воспользуемся следующими двумя утверждениями:

$F_{-k-1} = F_{-k+1} - F_{-k}$

$F_{-k} = F_{-k}$

Запишем в матричном виде:

$$\begin{pmatrix}
  F_{-k-1}\\ 
  F_{-k}
\end{pmatrix} = 
\begin{pmatrix}
  -1&1\\ 
  1&0
\end{pmatrix} *
\begin{pmatrix}
  F_{-k}\\ 
  F_{-k+1}
\end{pmatrix} =
{\begin{pmatrix}
  -1&1\\ 
  1&0
\end{pmatrix}}^2 *
\begin{pmatrix}
  F_{-k+1}\\ 
  F_{-k+2}
\end{pmatrix} = 
...=
{\begin{pmatrix}
  -1&1\\ 
  1&0
\end{pmatrix}}^k *
\begin{pmatrix}
  F_{-1}\\ 
  F_{0}
\end{pmatrix} = 
{\begin{pmatrix}
  -1&1\\ 
  1&0
\end{pmatrix}}^k *
\begin{pmatrix}
  1\\ 
  0
\end{pmatrix}
$$ 

Здесь у меня вычисляется $F_{-k}$ и $F_{-k-1}$ числа. если нужно именно $F_{-k}$ число, то используем

$$\begin{pmatrix}
  F_{-k}\\ 
  F_{-k+1}
\end{pmatrix} = 
{\begin{pmatrix}
  -1&1\\ 
  1&0
\end{pmatrix}}^{k-1} *
\begin{pmatrix}
  1\\ 
  0
\end{pmatrix}
$$

здесь на одну операцию меньше.

Если мы умеем умножать матрицы 2x2 за $O(1)$, а возводим матрицу в степень за $O(\log k)$, то $k$-е число анти-начи мы вычисляем за $O(\log k)$ (крч так же как и Фибоначчи, только матрица другая).

\bigskip

\section*{\textbf{Задание 2}} \textbf{\textit{Найдите асимптотическую оценку функции $T(n)$. Примените мастер-теорему в тех случаях, когда ее можно использовать, \textbf{и посчитайте асимптотику иначе, когда нельзя}. Варианты есть следующие. Можно выписать рекурренту в виде суммы и найти, чему она равна. Можно подставить рекурренту саму в себя и посмотреть, что получается. Можно обратиться к литературе (учебник Кормена, учебник Дасгупты)}}

Везде сначала проверяю мастер-теорему (выписываю a,b,c => проверяю условие на f(n) => делаю вывод => подсчитываю)

\medskip


\textbf{1}. $T(n) = 25 T(\frac{n}{5}) + n^2$

$a = 25, b = 5, c = \log_5 25 = 2$, тогда $ n^2 = \Theta (n^2) \Rightarrow T(n) = \Theta (n^{c} * \log n) = \Theta (n^{2} * \log n)$ по мастер-теореме

\bigskip

\textbf{2} $T(n) = 16 T(\frac{n}{2}) + n^3$

$a = 16, b = 2, c = \log_2 16 = 4$, тогда $ n^3 = O(n^{c-\varepsilon}) = O(n^{4 - 0.5}) \Rightarrow T(n) = \Theta (n^{c}) ) = \Theta (n^{4})$ по мастер-теореме (здесь подойдет любое $0 < \varepsilon < 1$, я написала 0.5 просто так)

\bigskip

\textbf{3} $T(n) = 9 T(\frac{n}{3}) + n^3$

$a = 9, b = 3, c = \log_3 9 = 2$, тогда $ n^3 = \Omega(n^{c+\varepsilon}) = \Omega(n^{2 + 0.5})$

В данном случае также нужно проверить условие регулярности: 
$k*f(n) \geq a*f(\frac{n}{b}) \Rightarrow k*n^3 \geq 9*(\frac{n}{3})^3 = \frac{n^3}{3} \Rightarrow $ выполняется при любом $1 > k > \frac{1}{3}$, тогда по условию мастер-теоремы 
$T(n) = \Theta (f(n)) = \Theta (n^{3})$

\bigskip


\textbf{4} $T(n) = T(n-1) + 3n$
В данном случае мы не можем использовать мастер-теорему, так как $T(n) \neq a*T(\frac{n}{b})$

Выпишем в виде суммы:
$T(n) = T(n-1) + 3n = T(n-2) + 3(n-1) + 3n = T(n-3) + 3(n-2) + 3(n-1) + 3n = ... = T(n-k) + 3(n-k+1) + 3(n-k+2) + ... + 3(n-1) + 3n = ... = 3*0 + 3*1 + 3*2 + ... + 3n = 3*(1 + 2 + ... + n) = 3 * \frac{n*(n+1)}{2} \Rightarrow T(n) = \Theta (n^2)$

\bigskip


\textbf{5} $T(n) = T(\frac{n}{4}) + T(\frac{3n}{4}) + n$

Докажем с помощью дерева рекурсии. Как отмечено в Кормене для аналогичного случаая, мы не можем утверждать, что наше дерево полное бинарное, а его листья расположены на разных уровнях: от $\log_4 n$ (самый близкий к корню, мы его получим по ветке n/4 -> n/4 -> n/4 ...) до $log_{4/3} n$ (самый далекий, его получим по ветке 3n/4 -> 3n/4 -> 3n/4 ...). Логика моего решения будет такая:

\begin{verbatim}
                   n
               /       \                -> n
           n/4          3n/4            -> n    # level 1
        /        \   /       \          -> n/4 + 3n/4 = n
     n/16     3n/16 3n/16    9n/16      -> n    # level 2
.......................................
 / \  / \ / \ / \ / \ / \ / \ / \ / \     # level log4 n
 O(1)  
.......................................
                                   / \     # level log4/3 n
                                O(1) O(1)
\end{verbatim}

Итак, а каждом уровне до $log_4 n$ проводится ровно n операций, который входят в f(n). На уровнях с $log_4 n$ до $log_{4/3} n$ листья выполняются за $O(1)$, а шагов из f(n) выполняется не более ($\leq n$), так как рекурсия из листьев не вызывается. Тогда мы можем оценить число шагов алгоритма как: $n \log_4 n  \leq T(n) \leq n \log_{4/3} n$, а, как мы доказывали раньше (типа в 1 дз что ли), это означает, что $T(n) = \Theta (n \log n)$ 

\bigskip


\textbf{6} $T(n) = T(\frac{n}{2}) + T(\frac{2n}{3}) + n^2$

Здесь придется немного посчитать. Будем решать с пмощью рекурретного дерева. Я НЕ ПИШУ $n^2$ на деревьях, потому что выглядело бы нечитаемо! 

\begin{verbatim}
                               n                    -> n^2
                           /       \                -> 
                       n/2          2n/3            -> n^2 * 7/6
                   /       \      /       \         -> 
                 n/4      n/3   n/3      4n/9       -> n^2 * (7/6)^2
               /    \    /   \  /  \    /    \      -> 
             n/8   n/6 n/6 2n/9 n/6 2n/9 2n/9  8n/27 -> n^2 * (7/6)^3
..............................................................................
 / \  / \ / \ / \ / \ / \ / \ / \ / \  / \  / \ / \      # level log2 n
 O(1)  
..............................................................................
                                                / \     # level log3/2 n
                                              O(1) O(1)
\end{verbatim}

Здесь мы заметим следующее. Сумма уровней в данном случае $\neq n$. Но мы все еще можем получить неплохую формулу :)

Давайте выпишем все слагаемые уровня k: их будет 2k, и каждое i и i+1 слагаемые (формально i четные, чтобы образовывалась пара) мы получаем из предыдущего уровня следующим способом:

$num_{k_i} = num_{k-1_i} * \frac{1}{2}; num_{k_{i+1}} = num_{k-1_i} * \frac{2}{3}$

Выпишем для всех слагаемых уровня $k$, рассматриваем попарно, здесь $i$ и $i+1$ элементы уровня k образуются из обшего предка на уровне k-1:

$$\sum_{i=0}^{k} num_{k-1_i} * \frac{1}{2} + num_{k-1_i} * \frac{2}{3} = \sum_{i=0}^{k}num_{k-1_i} * (\frac{1}{2} +\frac{2}{3}) = \sum_{i=0}^{k}\frac{7}{6} * num_{k-1_i} = \frac{7}{6} \sum_{i=0}^{k-1} num_{k-1_i} = ... $$ 

Но здесь мы брали k без ограничения общности, то есть справедливо, что аналогичная формула верна для уровне k-1, то тогда наша формула продолжится так (в конце $n^2$ = сумма 0 ряда = корня):

$$ ... = \frac{7}{6} \sum_{i=0}^{k-1} num_{k-1_i} = {\frac{7}{6}}^2 \sum_{i=0}^{k-2} num_{k-2_i} = ... = {\frac{7}{6}}^k n^2$$ 

Теперь мы можем считать сумму уровней, где нет листьев (от 0го до $\log_{2}n$), по формуле $\sum_k =  {\frac{7}{6}}^k n^2$. Тогда на уровнях, где есть листья (с $\log_{2}n$ по $\log_{3/2}n$, выполняется не более $\sum_k \leq  {\frac{7}{6}}^k n^2$ шагов. Тогда всего количество шагов:
$$\log_2 n * {\frac{7}{6}}^{\log_2 n} n^2 \leq T(n) \leq \log_{3/2} n * {\frac{7}{6}}^{\log_{3/2}} n^2$$. Здесь у нас встречается константа в степени $\log n$, мы не можем от нее избавиться. Поэтому я запишу $T(n) = \Theta (n^2\log n * {\frac{7}{6}}^{\log n})$ 

* нельзя написать $n^2 \log n$, потому что $n^2 \log n \leq n^2\log n * {\frac{7}{6}}^{\log n}$ для любых n. 

\bigskip


\textbf{7} $T(n) = T(n - 1) + n^2$

В данном случае мы не можем использовать мастер-теорему, так как $T(n) \neq a*T(\frac{n}{b})$

Выпишем в виде суммы:
$T(n) = T(n-1) + 3n^2 = T(n-2) + 3(n-1)^2 + 3n^2 = T(n-3) + 3(n-2)^2 + 3(n-1)^2 + 3n^2 = ... = T(n-k) + 3(n-k+1)^2 + 3(n-k+2)^2 + ... + 3(n-1)^2 + 3n^2 = ... = 3*0 + 3*1 + 3*2^2 + ... + 3n^2 = 3*(1 + 2^2 + ... + n^2) = 3 * \frac{n*(n+1)*(2n+1)}{6} \Rightarrow T(n) = \Theta (n^3)$

\bigskip


\textbf{8} $T(n) = 4 T(\frac{n}{16}) + \sqrt{n}$

$a = 4, b = 16, c = \log_{16} 4 = 0.5$, тогда $ \sqrt{n} = \Theta(n^{c}) = \Theta (n^{0.5}) = \Theta (\sqrt{n}) \Rightarrow T(n) = \Theta (n^{c} * \log n) = \Theta (n^{0.5} * \log n) = \Theta (\sqrt{n} * \log n)$


\bigskip


\textbf{9} $T(n) = 9 T(\frac{n}{3}) + n$

повтор п3

\bigskip


\section*{\textbf{Задание 3}} \textbf{\textit{Выведите первый случай мастер-теоремы (основной теоремы о рекуррентных соотношениях) для целых степеней (пренебрегая округлениями). Формально, нужно показать, что если}}

\begin{itemize}
    \item $T(n) = a T(\frac{n}{b}) + f(n)$, $1 \leq a \in \mathbb{N}$, $1 < b \in \mathbb{R}$
    \item $T(n) = \Theta(1)$ для $n < N_f$
    \item $\exists \varepsilon > 0: f(n) = O(n^{\log_b a - \varepsilon})$
\end{itemize}

то $T(n) = \Theta(n^{\log_b a})$

Ну начало доказательства мы вывели на лекции, поэтому совсем им не пользоваться было бы грешно :)

Сначала будем ставить рекурренту саму в себя (используем пункт 1):

$T(n) = a T(\frac{n}{b}) + f(n) = a^2 T(\frac{n}{b^2}) + a f(\frac{n}{b}) + f(n)= a^3 T(\frac{n}{b^3}) + a^2 f(\frac{n}{b^2}) + a f(\frac{n}{b}) + f(n) = ... = a^k T(\frac{n}{b^k}) + a^{k-1} f(\frac{n}{b^{k-1}}) +...+ a f(\frac{n}{b}) + f(n)$

Предположим, на k таком шаге мы дошли до значения $\frac{n}{b^k} = N_f$ из пункта 2, то есть $T(\frac{n}{b^k}) = T(N_f) = \Theta (1)$ согласно пункту 2. Теперь заменим везде $\frac{n}{b^k}$ на $N_f$, вместо $\Theta(1)$ в оценке я буду использовать $F$ в качестве обозначения константы (как было в лекции).

$a^k T(\frac{n}{b^k}) + a^{k-1} f(\frac{n}{b^{k-1}}) +...+ a f(\frac{n}{b}) + f(n) = a^k T(N_f) + a^{k-1} f(N_f b) + a^{k-2} f(N_f b^2) + ... + f(n) = a^kF + a^{k-1} f(N_f b) + a^{k-2} f(N_f b^2) + ... + f(n)$

Теперь мы можем воспользоваться пунктом 3: $\exists \varepsilon > 0: f(n) = O(n^{\log_b a - \varepsilon})$ и сделать оценку суммы (для каждого слагаемого верно, что $f(n) \leq P n^{c - \varepsilon}$, где P - положительная константа, а $c = \log_b a$. Без ограничения общности я буду использовать ее везде:

$a^kF + a^{k-1} f(N_f b) + a^{k-2} f(N_f b^2) + ... + f(n) \leq a^kF + a^{k-1} P{(N_f b)}^{c - \varepsilon} + a^{k-2} P{(N_f b^2)}^{c - \varepsilon} + ... + n^{c- \varepsilon}$

мы знаем, что $N_f = \frac{n}{b^k} \Rightarrow k = \log_b\frac{n}{N_f}$, тогда

$a^kF + a^{k-1} P{(N_f b)}^{c - \varepsilon} + a^{k-2} P{(N_f b^2)}^{c - \varepsilon} + ... + n^{c- \varepsilon} = a^{\log_b\frac{n}{N_f}}F + a^{\log_b\frac{n}{N_f}-1} P{(N_f b)}^{c - \varepsilon} + a^{\log_b\frac{n}{N_f}-2} P{(N_f b^2)}^{c - \varepsilon} + ... + n^{c- \varepsilon} = a^{\log_b n - \log_b N_f}F + a^{\log_b n - \log_b N_f - 1} P{(N_f b)}^{c - \varepsilon} + a^{\log_b n - \log_b N_f - 2} P{(N_f b^2)}^{c - \varepsilon} + ... + n^{c- \varepsilon} $

Теперь из каждого слагаемого, а точнее из коэффициента a в степени, мы можем вынести $a^{\log_b n}$ в качестве части, зависящей от n, а всю остальную часть запишем в константу F и будем ее теперь обозначать L (то есть L = P * a\_const\_part)

$a^{\log_b n}F + a^{\log_b n} L{(N_f b)}^{c - \varepsilon} + a^{\log_b n} L{(N_f b^2)}^{c - \varepsilon} + ... + n^{c- \varepsilon}$

Воспользуемся следующим свойством: $a^{\log_b n} = n^{\log_b a}$, снова обозначим $\log_b a = c$ и подставим в уравнение, а потом вспомним, что F = $\Theta(1)$.

$n^{\log_b a}F + n^{\log_b a} L{(N_f b)}^{c - \varepsilon} + n^{\log_b a} L{(N_f b^2)}^{c - \varepsilon} + ... + n^{c- \varepsilon} = Fn^c + Ln^c{(N_f b)}^{c - \varepsilon} + Ln^c{(N_f b^2)}^{c - \varepsilon} + ... + c^{c - \varepsilon} = \Theta(1)*n^c + Ln^c{(N_f b)}^{c - \varepsilon} + Ln^c{(N_f b^2)}^{c - \varepsilon} + ... + c^{c - \varepsilon} = \Theta(n^c) + L*n^c*({(N_f b^2)}^{c - \varepsilon} + {(N_f b^2)}^{c - \varepsilon} + ... + 1/L)$

Мы получили два слагаемых: $\Theta(n^c)$, что можно представить как $Const_1 n^c$, а во втором слагаемом за скобками находится $L n^c$, а внутри скобок константная часть, которую можно тоже обохначить какой-нить буквой, умножить на $L$, и тогда: 

$\Theta(n^c) + L*n^c*({(N_f b^2)}^{c - \varepsilon} + {(N_f b^2)}^{c - \varepsilon} + ... + 1/L) = \Theta(n^c) + H*O(n^c) =\Theta (n^c)$

вспомним, что $c = \log_b a$, и получим: 

$$T(n) = \Theta(n^{\log_b a})$$

\end{document}